\chapter{The Fitness Mode}\index{modes!Fitness Mode|textbf}
\label{ch:fitness}

\begin{versebox}
Every mode has ground to stand on,\\
every mode, a range to roam ---\\
test them all, and what holds up\\
is fit to take back home.
\end{versebox}

\noindent The first two modes showed you how to read: the Teaching Mode breaks a sutta into its six components; the Investigation Mode cross-examines questions and answers. But how do you know your reading is \textit{right}?\footnote{Pali: \textit{``Tattha katamo yuttihāro? `Sabbesaṃ hārānan'ti, ayaṃ yuttihāro.''} The mnemonic verse fragment is from Chapter~3: ``Every mode has ground and range'' (\textit{sabbesaṃ hārānaṃ yā bhūmi yo ca gocaro}). The word \textit{yutti} (fitness, applicability) comes from the root \textit{yuj}, ``to yoke, to connect.'' A reading ``fits'' when it can be yoked to the whole framework without contradiction.}

That is what the Fitness Mode is for. It is the reality check --- the quality control that every interpretation must pass before you can trust it. And the test is simple. Every reading must satisfy three criteria simultaneously.\footnote{Pali: \textit{``Kiṃ yojayati? Cattāro mahāpadesā buddhāpadeso saṅghāpadeso sambahulattherāpadeso ekattherāpadeso.''} The Netti names the Four Great Standards (\textit{cattāro mahāpadesā}) --- statements authorized by the Buddha, by the community, by a group of elders, or by a single elder --- and then immediately redefines them in terms of three tests. This passage closely parallels the famous Mahāparinibbāna Sutta (DN 16.4.8--11), where the Buddha establishes these same four standards as the criterion for authenticity after his passing.}

\section*{The Triple Test}

Any interpretation --- whether it comes from the Buddha himself, from the community, from a group of elders, or from a single elder --- must be checked against three things.\footnote{Pali: \textit{``Ime cattāro mahāpadesā, tāni padabyañjanāni sutte otārayitabbāni, vinaye sandassayitabbāni, dhammatāyaṃ upanikkhipitabbāni.''} Literally: ``These four great standards --- those words and phrases must be \textit{brought down into} the suttas, \textit{displayed alongside} the discipline, and \textit{placed within} the nature of things.'' Three verbs, three tests, three domains.}

First: does it fit the suttas? Specifically, does it land within the Four Noble Truths\index{Four Noble Truths|textbf}?\footnote{Pali: \textit{``Katamasmiṃ sutte otārayitabbāni? Catūsu ariyasaccesu.''} The ``sutta'' here is not a specific text but the doctrinal framework itself. Any authentic teaching must be traceable to suffering, its origin, its cessation, or the path.}

Second: does it fit the discipline? Meaning: does it lead to the removal of greed, hatred, and delusion\index{greed}\index{hatred}\index{delusion}?\footnote{Pali: \textit{``Katamasmiṃ vinaye sandassayitabbāni? Rāgavinaye dosavinaye mohavinaye.''} The Netti reads \textit{vinaya} not as the monastic code but as the ``removal'' (\textit{vinaya} from \textit{vi-nī}) of the three roots of unwholesomeness. This is an important hermeneutical move: ``discipline'' means whatever removes defilements.}

Third: does it fit the nature of things? That is: does it align with dependent origination\index{dependent origination|textbf}?\footnote{Pali: \textit{``Katamissaṃ dhammatāyaṃ upanikkhipitabbāni? Paṭiccasamuppāde.''} The ``nature of things'' (\textit{dhammatā}) is equated directly with dependent origination (\textit{paṭiccasamuppāda}). Any reading that contradicts conditionality fails the fitness test.}

If a reading descends into the Four Noble Truths, if it is visible in the removal of defilements, and if it does not contradict the nature of things --- then it does not generate further corruptions. Whatever fits by means of the Four Great Standards, in whatever way it fits, by whatever reasoning it fits --- that should be accepted.\footnote{Pali: \textit{``Yadi catūsu ariyasaccesu avatarati, kilesavinaye sandissati, dhammatañca na vilometi, evaṃ āsave na janeti. Catūhi mahāpadesehi yaṃ yaṃ yujjati, yena yena yujjati, yathā yathā yujjati, taṃ taṃ gahetabbaṃ.''} The triple \textit{yaṃ yaṃ \ldots yena yena \ldots yathā yathā} construction (``whatever \ldots by whatever \ldots in whatever way'') emphasizes the universality of the test. There is no fourth criterion. Three is enough.}

\section*{Same Words, Same Thing?}

Now the Fitness Mode turns to a subtler question: when a sutta uses multiple terms, how do you know whether they mean the same thing or different things?\footnote{Pali: \textit{``Pañhaṃ pucchitena kati padāni pañheti padaso pariyogāhitabbaṃ vicetabbaṃ?''} ``When a question is asked, how many terms are in the question? This should be examined term by term.'' The Netti is setting up its method for semantic analysis: are multiple terms pointing at one meaning, or do they each carry a distinct meaning?}

If all the terms in a question point to a single meaning, it is one question. If four terms point to one meaning, one question. Three terms, one meaning --- one question. Two terms --- one question. One term, one meaning --- one question. The reader must examine: do these terms have different meanings \textit{and} different wording? Or do they share one meaning with only the wording being different?\footnote{Pali: \textit{``Taṃ upaparikkhamānena aññātabbaṃ kiṃ ime dhammā nānatthā nānābyañjanā, udāhu imesaṃ dhammānaṃ eko attho byañjanameva nānan'ti.''} This is the key diagnostic question of the Fitness Mode: \textit{nānatthā nānābyañjanā} (different in meaning and different in phrasing) versus \textit{eko attho byañjanameva nānaṃ} (one meaning, only the phrasing is different). Getting this distinction right is essential for accurate interpretation.}

Here is an example. A deity asks the Buddha:\footnote{This verse is from SN 1.66 (\textit{Abbhāhata Sutta}). The Netti regularly draws on the Saṃyutta Nikāya's Devatā Saṃyutta for its worked examples --- these short, punchy deity-verses are ideal for demonstrating hermeneutical technique because their compression forces close reading.}

\begin{versebox}
``By what is the world struck down?\\
By what is it hemmed in?\\
What is the dart that pierces it?\\
What keeps it fuming without end?''
\end{versebox}

Four terms were asked. But actually, this is three questions --- and the way you know is from the Buddha's answer:\footnote{Pali: \textit{``Imāni cattāri padāni pucchitāni. Te tayo pañhā kathaṃ ñāyati? Bhagavā hi devatāya visajjeti.''} ``These four terms were asked. How is it known that they are three questions? Because the Blessed One answers the deity.'' The Netti's reasoning: the question has four lines, but the answer reveals that two of the four lines are actually about the same thing.}

\begin{versebox}
``By death the world is struck down.\\
By aging it is hemmed in.\\
Pierced by the dart of craving,\\
it fumes with wanting without end.''
\end{versebox}

\section*{Aging Is Not Death, and Neither Is Craving}

Now the Fitness Mode demonstrates its diagnostic power. Aging and death --- these are two characteristics of conditioned things. Aging is the alteration of what persists; death is the ending. But they differ in meaning.\footnote{Pali: \textit{``Tattha jarā ca maraṇañca imāni dve saṅkhatassa saṅkhatalakkhaṇāni. Jarāyaṃ ṭhitassa aññathattaṃ, maraṇaṃ vayo.''} Two of the three marks of the conditioned (\textit{saṅkhatalakkhaṇa}): change-while-standing (\textit{ṭhitassa aññathattaṃ}) is aging; destruction (\textit{vaya}) is death. The third mark --- arising (\textit{uppāda}) --- is not mentioned here because it is not relevant to the verse being analyzed.}

How so? Even those still in the womb can die, though they have not aged. The gods die, but their bodies do not grow old. You can sometimes reverse aging, but you cannot reverse death --- except, the Netti adds drily, within the domain of those who possess psychic powers.\footnote{Pali: \textit{``Gabbhagatāpi hi mīyanti, na ca te jiṇṇā bhavanti. Atthi ca devānaṃ maraṇaṃ, na ca tesaṃ sarīrāni jīranti. Sakkateva jarāya paṭikammaṃ kātuṃ, na pana sakkate maraṇassa paṭikammaṃ kātuṃ aññatreva iddhimantānaṃ iddhivisayā.''} The phrase \textit{aññatreva iddhimantānaṃ iddhivisayā} (``except within the domain of the psychic-power-possessors'') is a characteristically wry aside. The Netti's logical method is visible here: aging $\neq$ death because each can occur without the other. This is a classic dissociation argument (\textit{vinibbhoga}).}

And what about the dart of craving? That is different from both aging and death. The proof: there are those free from craving who still grow old and die. If craving were the same as aging and death, then everyone who lost their youth would also be free of craving. And if aging and death were the origin of suffering --- as craving is --- then craving would not be the origin of suffering. But craving \textit{is} the origin of suffering, not aging and death. And if aging and death could be destroyed by the path --- as craving can --- then they too would be path-destroyable. By this kind of fitness-reasoning, one must search through mutual reasons.\footnote{Pali: \textit{``Yaṃ panāha taṇhāsallena otiṇṇo'ti dissanti vītarāgā jīrantāpi mīyantāpi. Yadi ca yathā jarāmaraṇaṃ, evaṃ taṇhāpi siyā. Evaṃ sante sabbe yobbanaṭṭhāpi vigatataṇhā siyuṃ. Yathā ca taṇhā dukkhassa samudayo \ldots na hi jarāmaraṇaṃ dukkhassa samudayo, taṇhā dukkhassa samudayo. Yathā ca taṇhā maggavajjhā, evaṃ jarāmaraṇampi siyā maggavajjhaṃ. Imāya yuttiyā aññamaññehi kāraṇehi gavesitabbaṃ.''} This is the Fitness Mode at work: a chain of logical arguments (\textit{yutti}) showing that craving is categorically different from aging and death. The reasoning is precise: (1) arahants age and die but have no craving, so craving $\neq$ aging-and-death; (2) craving is the origin of suffering (Second Noble Truth), but aging-and-death is not; (3) craving is destroyable by the path, but aging-and-death is not (arahants still age and die). Three arguments, three angles, one conclusion.}

If the fitness-reasoning confirms that there is a difference in meaning, then one should also search the wording.\footnote{Pali: \textit{``Yadi ca sandissati yuttisamārūḷhaṃ atthato ca aññattaṃ, byañjanatopi gavesitabbaṃ.''} Having established difference of meaning, the Netti says to also check the phrasing --- a two-step method. First test the semantics, then confirm with the language.}

\section*{The Dart and the Fuming}

But ``dart'' and ``fuming'' --- these two terms share a single meaning. It does not fit to say that wanting and craving differ in meaning.\footnote{Pali: \textit{``Salloti vā dhūpāyananti vā imesaṃ dhammānaṃ atthato ekattaṃ. Na hi yujjati icchāya ca taṇhāya ca atthato aññattaṃ.''} The Netti's diagnosis: ``dart'' (\textit{salla} = craving) and ``fuming'' (\textit{dhūpāyana} = wanting/desire) are two names for the same thing. The four-line question turns out to contain three distinct topics: (1) death, (2) aging, and (3) craving-wanting. Hence: four terms, three questions.}

When craving's purpose is unfulfilled, anger and resentment arise at nine occasions for annoyance.\footnote{Pali: \textit{``Taṇhāya adhippāye aparipūramāne navasu āghātavatthūsu kodho ca upanāho ca uppajjati.''} The nine occasions for annoyance (\textit{nava āghātavatthūni}) are a standard list: ``He harmed me, he harms me, he will harm me; he harmed someone dear to me, he harms someone dear to me, he will harm someone dear to me; he helped someone I dislike, he helps someone I dislike, he will help someone I dislike.'' The Netti's point: even anger is ultimately rooted in frustrated craving. By this fitness-reasoning, aging and death differ from craving in meaning.} By this fitness-reasoning, aging, death, and craving differ from each other in meaning.

Now the Buddha used two names for what is essentially one thing --- ``wanting'' and ``craving.'' He did this because of the different external objects that serve as their basis.\footnote{Pali: \textit{``Idaṃ bhagavatā bāhirānaṃ vatthūnaṃ ārammaṇavasena dvīhi nāmehi abhilapitaṃ icchātipi taṇhātipi.''} The Netti explains the semantic relationship: ``wanting'' (\textit{icchā}) and ``craving'' (\textit{taṇhā}) are not two different mental states but one state with two names, each name emphasizing a different aspect of its object.} But all craving shares a single defining characteristic: clinging.

\section*{Like Fire, Like Craving}

Think of fire. All fire shares one characteristic: heat. But depending on what it burns, it gets different names --- wood-fire, grass-fire, kindling-fire, dung-fire, chaff-fire, rubbish-fire. All fire is simply heat.\footnote{Pali: \textit{``Yathā sabbo aggi uṇhattalakkhaṇena ekalakkhaṇo, api ca upādānavasena aññamaññāni nāmāni labhati, kaṭṭhaggītipi tiṇaggītipi sakalikaggītipi gomayaggītipi thusaggītipi saṅkāraggītipi, sabbo hi aggi uṇhattalakkhaṇo'va.''} The fire analogy is one of the Netti's most memorable illustrations. The word for ``fuel'' (\textit{upādāna}) is the same word used for ``clinging'' in the dependent origination formula --- a double meaning the Netti is certainly aware of. Just as fire is named by its fuel, craving is named by its object.}

In the same way, all craving shares one characteristic --- clinging --- but depending on its object it gets called by different names: wanting, craving, the dart, the fuming, the current, the sticky trap, affection, exhaustion, the creeper, conceit, the bond, hope, thirst, delight.\footnote{Pali: \textit{``Evaṃ sabbā taṇhā ajjhosānalakkhaṇena ekalakkhaṇā, api tu ārammaṇaupādānavasena aññamaññehi nāmehi abhilapitā icchāitipi taṇhāitipi salloitipi dhūpāyanāitipi saritāitipi visattikāitipi sinehoitipi kilamathoitipi latāitipi maññanāitipi bandhoitipi āsāitipi pipāsāitipi abhinandanāitipi.''} Fourteen synonyms for craving. Several of these are botanical metaphors (\textit{latā} = creeper, vine), others physiological (\textit{pipāsā} = thirst, \textit{sinehā} = moisture/affection), others experiential (\textit{kilamatha} = exhaustion, \textit{maññanā} = conceit). The diversity of names reflects the diversity of objects; the single characteristic reflects the unity of the underlying state.}

All craving shares one characteristic: clinging. As the synonyms verse puts it:\footnote{This verse is widely quoted in the Pali Canon; cf.\ Dhp-a and various Nikāya passages. The Netti cites it as \textit{vevacana} (synonym-verse) evidence that the many names for craving all point to one thing.}

\begin{versebox}
Hope, longing, and delight ---\\
currents lodged in many realms ---\\
born from the root of ignorance,\\
all yearnings I have ended, root and stem.
\end{versebox}

These are all synonyms for craving. As the Buddha said: ``Tissa, regarding form --- for one who has not abandoned lust, who has not abandoned desire, affection, thirst, and fever \ldots'' --- and the same for feeling, perception, formations, and consciousness. The entire sutta should be expanded thus. These are all synonyms for craving. This is how fitness works.\footnote{Pali: \textit{``Taṇhāyetaṃ vevacanaṃ. Yathāha bhagavā --- rūpe tissa avigatarāgassa avigatacchandassa avigatapemassa avigatapipāsassa avigatapariḷāhassa \ldots Taṇhāyetaṃ vevacanaṃ. Evaṃ yujjati.''} The Netti demonstrates: when the Buddha lists five terms --- lust (\textit{rāga}), desire (\textit{chanda}), affection (\textit{pema}), thirst (\textit{pipāsā}), fever (\textit{pariḷāha}) --- for each of the five aggregates, these are not five different mental states. They are five synonyms for craving, applied to five different objects. The Fitness Mode helps you see that 5 $\times$ 5 = 25 apparent distinctions are really one thing named in twenty-five ways.}

\section*{The Right Medicine for the Right Patient}

All the approach to suffering is rooted in sensual craving and its formations. But it does not fit to say that all the approach to disenchantment is likewise rooted in sensual craving and its supports. By this fitness-reasoning, one must search through mutual reasons.\footnote{Pali: \textit{``Sabbo dukkhūpacāro kāmataṇhāsaṅkhāramūlako, na pana yujjati sabbo nibbidūpacāro kāmataṇhāparikkhāramūlako.''} A subtle point: while all suffering flows from craving, not all disenchantment flows from the same source. Different people become disillusioned through different gates. This asymmetry is important --- it means the path out is not simply the mirror image of the path in.}

The Buddha teaches contemplation of the unattractive to a person dominated by greed. He teaches loving-kindness to a person dominated by hatred\index{loving-kindness}. He teaches dependent origination to a person dominated by delusion.\footnote{Pali: \textit{``Yathā hi bhagavā rāgacaritassa puggalassa asubhaṃ desayati, dosacaritassa bhagavā puggalassa mettaṃ desayati. Mohacaritassa bhagavā puggalassa paṭiccasamuppādaṃ desayati.''} The three character types (\textit{carita}) and their remedies are a fundamental principle of the Pali teaching on meditation: greed-types get the unattractive (\textit{asubha}), hate-types get loving-kindness (\textit{mettā}), delusion-types get dependent origination (\textit{paṭiccasamuppāda}). See Vism III for the full treatment.}

Suppose the Buddha were to teach the liberation of the heart through loving-kindness to a greed-dominated person, or the pleasant path with slow insight, or the pleasant path with quick insight, or abandonment preceded by insight-meditation. It would not fit. Whatever counteracts greed, whatever counteracts hatred, whatever counteracts delusion --- all of that should first be investigated with the Investigation Mode, then tested for fitness with the Fitness Mode. That is the extent of wisdom's domain.\footnote{Pali: \textit{``Yadi hi bhagavā rāgacaritassa puggalassa mettaṃ cetovimuttiṃ deseyya \ldots na yujjati desanā. Evaṃ yaṃ kiñci rāgassa anulomappahānaṃ dosassa anulomappahānaṃ mohassa anulomappahānaṃ. Sabbaṃ taṃ vicayena hārena vicinitvā yuttihārena yojetabbaṃ. Yāvatikā ñāṇassa bhūmi.''} The Netti gives us the methodology: (1) investigate with \textit{vicaya}, (2) test with \textit{yutti}. And the closing phrase --- \textit{yāvatikā ñāṇassa bhūmi}, ``that is the extent of wisdom's domain'' --- quietly claims that this two-step method covers all of doctrinal understanding.}

\section*{What Fits and What Doesn't}

Now the Netti does something clever. It takes the four ``divine abodes'' --- loving-kindness, compassion, sympathetic joy, and equanimity --- and shows that each one is a targeted antidote to one specific poison. Get the pairing wrong, and the medicine doesn't work.\footnote{The original Pali for this entire passage (\S21) follows a strict binary pattern: \textit{na yujjati desanā} (``the teaching does not fit'') / \textit{yujjati desanā} (``the teaching fits''), repeated for each case. We have recast this repetitive litany as a narrative summary because the pattern, once grasped, does not need to be rehearsed twelve times in English. The full Pali is preserved in the notes for those who want the original cadence. The gist is identical.}

Loving-kindness specifically dissolves ill-will. If you are practicing loving-kindness, it would be absurd to claim that ill-will is gaining ground in your mind --- that does not fit. What \textit{does} fit is that ill-will is being abandoned. Compassion dissolves cruelty. Sympathetic joy dissolves boredom and discontent. Equanimity dissolves lust. Each pairing is precise; swap them around and the diagnosis fails.\footnote{Pali: \textit{``Mettāvihārissa sato byāpādo cittaṃ pariyādāya ṭhassatīti na yujjati desanā, byāpādo pahānaṃ abbhatthaṃ gacchatīti yujjati desanā. Karuṇāvihārissa sato vihesā \ldots Muditāvihārissa sato arati \ldots Upekkhāvihārissa sato rāgo \ldots''} The four divine abodes (\textit{brahmavihāra}) and their specific enemies: loving-kindness (\textit{mettā}) vs.\ ill-will (\textit{byāpāda}); compassion (\textit{karuṇā}) vs.\ cruelty (\textit{vihesā}); sympathetic joy (\textit{muditā}) vs.\ discontent (\textit{arati}); equanimity (\textit{upekkhā}) vs.\ lust (\textit{rāga}).}

The same logic applies to the more advanced practices. Someone dwelling in the signless meditation --- a state where the mind has dropped all mental ``signs'' and labels --- is not going to be overwhelmed by sign-chasing. And someone who has genuinely let go of the sense ``I am'' is not going to be overwhelmed by doubt about who or what they are.\footnote{Pali: \textit{``Animittavihārissa sato nimittānusārī tena teneva viññāṇaṃ pavattatīti na yujjati desanā \ldots Asmīti vigataṃ ayamahamasmīti na samanupassāmi. Atha ca pana me kismīti kathasmīti vicikicchā kathaṃkathāsallaṃ cittaṃ pariyādāya ṭhassatīti na yujjati desanā, vicikicchā kathaṃkathāsallaṃ pahānaṃ abbhatthaṃ gacchatīti yujjati desanā.''} The signless meditation (\textit{animitta vihāra}) is a high-level practice in which the meditator drops all mental ``signs'' (\textit{nimitta}) or proliferations. The case about identity-view and doubt is especially interesting: once the conceit ``I am'' (\textit{asmi-māna}) is gone, the question ``What am I?'' simply cannot arise.}

The point is not that these pairings are obvious --- some of them are --- but that the Fitness Mode gives you a \textit{method} for checking. You don't need to guess whether a teaching fits a situation. You test it.

\section*{The Ladder of Absorption}

Now the Netti applies the same fitness test to the entire ladder of meditative absorption\index{jhāna|textbf} --- and here the pattern becomes beautifully clear.

Think of a staircase. At each step, what is \textit{below} you can only pull you down. What is \textit{above} you is the only direction worth moving. To claim otherwise --- that falling backward would lead you higher --- simply does not fit.\footnote{Pali: \textit{``Yathā vā pana paṭhamaṃ jhānaṃ samāpannassa sato kāmarāgabyāpādā visesāya saṃvattantīti na yujjati desanā, hānāya saṃvattantīti yujjati desanā. Vitakkasahagatā vā saññāmanasikārā hānāya saṃvattantīti na yujjati desanā, visesāya saṃvattantīti yujjati desanā.''} The original Pali applies this pattern individually to each of the eight attainments (four form absorptions and four formless bases), spelling out each case in full. The repetition is meditative in its own way, but in English it reads more naturally as a summary of the principle, with the specifics in the notes. \textbf{A note on our approach}: the original is considerably more repetitive than what we present here --- it works through all eight levels one by one in the same binary ``fits / does not fit'' formula. We have distilled the principle and given examples from the bottom and top of the ladder. Serious students should consult the Pali for the complete iteration, which has a mantra-like quality that rewards slow reading.}

At the first absorption, sensual desire and ill-will belong to the level below --- they can only cause decline. But thought and examination, which are still present in the first absorption, belong at this level and point upward. At the second absorption, it is thought and examination that now belong below --- they cause decline. But equanimity and happiness point upward. And so on, rung by rung, through all four form absorptions and all four formless attainments.\footnote{The full sequence in the Pali: (1) First absorption: sensual desire/ill-will → decline; thought-accompanied perceptions → distinction. (2) Second absorption: thought-accompanied perceptions → decline; equanimity-happiness perceptions → distinction. (3) Third absorption: rapture-accompanied perceptions → decline; equanimity-purity-of-mindfulness → distinction. (4) Fourth absorption: equanimity-accompanied perceptions → decline; perceptions of the base of infinite space → distinction. (5) Base of infinite space: form-accompanied perceptions → decline; perceptions of infinite consciousness → distinction. (6) Base of infinite consciousness: perceptions of infinite space → decline; perceptions of nothingness → distinction. (7) Base of nothingness: perceptions of infinite consciousness → decline; perceptions of neither-perception-nor-non-perception → distinction. (8) Base of neither-perception-nor-non-perception: perception-activity → decline; perceptions of the cessation of perception and feeling → distinction.}

At the very top of the ladder --- the base of neither-perception-nor-non-perception, the most refined state of consciousness possible --- even the faintest flicker of perception-activity would mean going backward. The only thing that leads higher from there is the cessation of perception and feeling itself: the mind's final, voluntary silence.\footnote{Pali: \textit{``Nevasaññānāsaññāyatanaṃ samāpannassa sato saññūpacārā visesāya saṃvattantīti na yujjati desanā \ldots Saññāvedayitanirodhasahagatā vā saññāmanasikārā hānāya saṃvattantīti na yujjati desanā, visesāya saṃvattantīti yujjati desanā.''} The cessation of perception and feeling (\textit{saññāvedayitanirodha}) is the supreme meditative achievement, attainable only by non-returners and arahants. It borders on nirvana itself.}

And when the mind has been thoroughly trained through this entire progression --- when it is ripe, supple, ready --- then the great breakthrough fits. A mind that has climbed the whole ladder is \textit{naturally} ready for liberation. To say otherwise would be like saying a fully ripe fruit is not ready to fall.\footnote{Pali: \textit{``Kallatāparicitaṃ cittaṃ na ca abhinīhāraṃ khamatīti na yujjati desanā, kallatāparicitaṃ cittaṃ atha ca abhinīhāraṃ khamatīti yujjati desanā.''} The ``great breakthrough'' (\textit{abhinīhāra}) refers to directing the concentrated mind toward liberating insight and the supramundane path. The Pali word \textit{kallatāparicitaṃ} (``trained to readiness, exercised to pliancy'') conveys a mind that has been worked, shaped, tempered --- like metal heated and hammered until it will hold any form.}

\ornament

In this way, all nine suttas should be --- according to the teaching, according to the discipline, according to the Teacher's instruction --- thoroughly investigated with the Investigation Mode, then tested for fitness with the Fitness Mode. That is why the Venerable Mahākaccāyana said: ``Every mode has ground to stand on, every mode a range to roam.''\footnote{Pali: \textit{``Evaṃ sabbe navasuttantā yathādhammaṃ yathāvinayaṃ yathāsatthusāsanaṃ sabbato vicayena hārena vicinitvā yuttihārena yojetabbāti. Tenāha āyasmā mahākaccāyano `sabbesaṃ hārānaṃ yā bhūmi yo ca gocaro tesan'ti.''} The closing formula ties back to the mnemonic verse and establishes the Fitness Mode's relationship to the Investigation Mode: first investigate, then test for fitness. The two modes work as a pair --- analysis followed by validation.}
